% ========================================
% 简化版教材LaTeX模板 - 确保编译成功
% 版本: v1.0
% 适用: 大学教材编写 (简化版)
% 更新: 2025年8月28日
% ========================================

% 文档类配置 - 中文书籍类
\documentclass[
    12pt,           % 字体大小
    a4paper,        % 纸张大小
    twoside,        % 双面排版
    openright       % 章节从右页开始
]{ctexbook}

% ========================================
% 基础包引入 - 仅包含必需包
% ========================================

% 数学支持
\usepackage{amsmath,amssymb,amsthm}

% 图形和颜色
\usepackage{graphicx}
\usepackage{xcolor}

% 页面布局
\usepackage[
    top=2.5cm,
    bottom=2.5cm,
    left=2.8cm,
    right=2.2cm,
    bindingoffset=0.5cm,
    headheight=15pt,
    footskip=1.5cm
]{geometry}

% 字体和编码
\usepackage{fontspec}
\usepackage{setspace}
\onehalfspacing  % 1.5倍行距

% 表格增强
\usepackage{longtable}
\usepackage{booktabs}
\usepackage{array}
\usepackage{calc}

% 列表环境
\usepackage{enumitem}

% 代码环境
\usepackage{listings}

% 页眉页脚
\usepackage{fancyhdr}

% 超链接
\usepackage{hyperref}

% ========================================
% 中文字体配置 - 简化版
% ========================================

% 设置中文字体（简化配置）
\setCJKmainfont{SimSun}[
    BoldFont=SimHei,
    ItalicFont=KaiTi
]
\setCJKsansfont{SimHei}
\setCJKmonofont{FangSong}

% 英文字体配置
\setmainfont{Times New Roman}
\setsansfont{Arial}
\setmonofont{Courier New}

% ========================================
% 配色方案定义
% ========================================

% 主色调
\definecolor{primarycolor}{RGB}{52, 152, 219}      % 主蓝色
\definecolor{secondarycolor}{RGB}{230, 126, 34}    % 橙色
\definecolor{accentcolor}{RGB}{39, 174, 96}        % 绿色
\definecolor{warningcolor}{RGB}{241, 196, 15}      % 黄色
\definecolor{dangercolor}{RGB}{231, 76, 60}        % 红色

% 灰度系列
\definecolor{lightgray}{RGB}{248, 249, 250}
\definecolor{mediumgray}{RGB}{173, 181, 189}
\definecolor{darkgray}{RGB}{73, 80, 87}

% 代码配色
\definecolor{codeblue}{RGB}{40, 90, 200}
\definecolor{codegreen}{RGB}{0, 150, 0}
\definecolor{codered}{RGB}{200, 0, 0}
\definecolor{codegray}{RGB}{128, 128, 128}
\definecolor{backcolour}{RGB}{248, 248, 248}

% ========================================
% 页面样式配置
% ========================================

% 页眉页脚样式
\pagestyle{fancy}
\fancyhf{}
\fancyhead[LE]{\leftmark}
\fancyhead[RO]{\rightmark}
\fancyfoot[LE,RO]{\thepage}
\renewcommand{\headrulewidth}{0.4pt}
\renewcommand{\footrulewidth}{0pt}

% 章节样式
\ctexset{
    chapter={
        format={\centering\Huge\bfseries},
        name={第,章},
        number={\arabic{chapter}},  % 使用阿拉伯数字
        beforeskip=20pt,
        afterskip=40pt
    },
    section={
        format={\Large\bfseries},
        beforeskip=20pt,
        afterskip=10pt
    },
    subsection={
        format={\large\bfseries},
        beforeskip=15pt,
        afterskip=8pt
    }
}

% ========================================
% 定理环境配置
% ========================================

% 定义定理样式
\theoremstyle{definition}
\newtheorem{definition}{定义}[chapter]
\newtheorem{theorem}{定理}[chapter]
\newtheorem{lemma}{引理}[chapter]
\newtheorem{example}{例}[chapter]
\newtheorem{exercise}{练习}[chapter]

% ========================================
% 列表样式配置
% ========================================

% 设置列表样式
\setlist[itemize,1]{label=\textbullet}
\setlist[itemize,2]{label=\textendash}
\setlist[itemize,3]{label=\textasteriskcentered}

\setlist[enumerate,1]{label=\arabic*.}
\setlist[enumerate,2]{label=(\alph*)}
\setlist[enumerate,3]{label=\roman*.}

% ========================================
% 自定义命令 - 简化版
% ========================================

% Pandoc兼容性命令
\providecommand{\tightlist}{%
  \setlength{\itemsep}{0pt}\setlength{\parskip}{0pt}}

% 强调命令（简化版，不依赖特殊字体）
\newcommand{\highlight}[1]{\textcolor{primarycolor}{\textbf{#1}}}
\newcommand{\important}[1]{\textcolor{dangercolor}{\textbf{#1}}}
\newcommand{\note}[1]{\textcolor{secondarycolor}{\textit{#1}}}

% 代码相关命令
\newcommand{\code}[1]{\texttt{#1}}
\newcommand{\file}[1]{\texttt{\textbf{#1}}}
\newcommand{\variable}[1]{\texttt{\textit{#1}}}

% ========================================
% 简化的特殊环境定义
% ========================================

% 学习目标（简化版 - 不使用tcolorbox）
\newcommand{\learningobjectives}[1]{
    \begin{quote}
        \textcolor{primarycolor}{\textbf{【学习目标】}}
        \par
        #1
    \end{quote}
}

% 关键概念（简化版）
\newcommand{\keypoints}[1]{
    \begin{quote}
        \textcolor{secondarycolor}{\textbf{【关键概念】}}
        \par
        #1
    \end{quote}
}

% 实践练习（简化版）
\newcommand{\practiceexercise}[1]{
    \begin{quote}
        \textcolor{accentcolor}{\textbf{【实践练习】}}
        \par
        #1
    \end{quote}
}

% 注意事项（简化版）
\newcommand{\attention}[1]{
    \begin{quote}
        \textcolor{warningcolor}{\textbf{【注意】}}
        \par
        #1
    \end{quote}
}

% 章节摘要（简化版）
\newcommand{\chaptersummary}[1]{
    \begin{quote}
        \textcolor{primarycolor}{\textbf{【本章要点】}}
        \par
        #1
    \end{quote}
}

% ========================================
% 代码高亮配置 - 简化版
% ========================================

\lstset{
    backgroundcolor=\color{backcolour},
    commentstyle=\color{codegreen}\itshape,
    keywordstyle=\color{codeblue}\bfseries,
    numberstyle=\tiny\color{codegray},
    stringstyle=\color{codered},
    basicstyle=\ttfamily\scriptsize,       % 使用较小字体
    breakatwhitespace=false,
    breaklines=true,
    captionpos=b,
    keepspaces=true,
    numbers=left,
    numbersep=8pt,
    showspaces=false,
    showstringspaces=false,
    showtabs=false,
    tabsize=4,
    frame=single,
    rulecolor=\color{mediumgray},
    framexleftmargin=8pt,
    framexrightmargin=8pt,
    framextopmargin=8pt,
    framexbottommargin=8pt,
    xleftmargin=15pt,
    xrightmargin=15pt
}

% 基础代码语言样式
\lstdefinestyle{csharp}{
    language=[Sharp]C,
    keywordstyle=\color{codeblue}\bfseries,
    commentstyle=\color{codegreen}\itshape,
    stringstyle=\color{codered}
}

\lstdefinestyle{javascript}{
    language=JavaScript,
    keywordstyle=\color{codeblue}\bfseries,
    commentstyle=\color{codegreen}\itshape,
    stringstyle=\color{codered}
}

\lstdefinestyle{python}{
    language=Python,
    keywordstyle=\color{codeblue}\bfseries,
    commentstyle=\color{codegreen}\itshape,
    stringstyle=\color{codered}
}

\lstdefinestyle{java}{
    language=Java,
    keywordstyle=\color{codeblue}\bfseries,
    commentstyle=\color{codegreen}\itshape,
    stringstyle=\color{codered}
}

\lstdefinestyle{sql}{
    language=SQL,
    keywordstyle=\color{codeblue}\bfseries,
    commentstyle=\color{codegreen}\itshape,
    stringstyle=\color{codered}
}

\lstdefinestyle{terminal}{
    basicstyle=\ttfamily\scriptsize\color{white},
    backgroundcolor=\color{black},
    showstringspaces=false,
    numbers=none,
    frame=single
}

% ========================================
% 超链接配置
% ========================================

\hypersetup{
    colorlinks=true,
    linkcolor=primarycolor,
    filecolor=primarycolor,
    urlcolor=primarycolor,
    citecolor=primarycolor,
    bookmarksnumbered=true,
    bookmarksopen=true
}

% ========================================
% 文档配置
% ========================================

% 设置段落缩进
\setlength{\parindent}{2em}
\setlength{\parskip}{0.5ex plus 0.2ex minus 0.1ex}

% 防止过长行
\setlength{\emergencystretch}{3em}
\setcounter{secnumdepth}{5}

% 图表标题样式
\usepackage{caption}
\captionsetup{
    font=small,
    labelfont=bf,
    textfont=it,
    margin=10pt
}

% ========================================
% 模板使用说明
% ========================================

% 本模板是简化版，移除了可能导致编译问题的包：
% - 移除了 minted 包（需要 -shell-escape）
% - 移除了 fontawesome5 包（可能在某些系统中不可用）
% - 移除了 tcolorbox 包（简化了特殊环境）
% - 移除了 tikz 包（减少依赖）

% 可用的特殊环境：
% - \learningobjectives{内容} - 学习目标框
% - \keypoints{内容} - 关键概念框  
% - \practiceexercise{内容} - 实践练习框
% - \attention{内容} - 注意事项框
% - \chaptersummary{内容} - 章节摘要框

% 常用命令：
% - \highlight{文本} - 高亮强调
% - \important{文本} - 重要内容
% - \note{文本} - 注释说明
% - \code{代码} - 行内代码